\section{Statements and conventions}
Write $\gamma_N$ for standard Gaussian measure on $\R^N$ and $\cP_2(\R^N)$
for probability laws with finite second moment. For $\nu\in\cP_2(\R^N)$ set
\[
 \mathcal F_\nu(C)=\sup\{\E[X\cdot Y]:X\sim N(0,C),\ Y\sim\nu\},
 \qquad C\geq0.
\]
The supremum is over all couplings. Covariance matrices are varied in their
ordinary linear coordinates.

\begin{theorem}[Covariance concavity]\label{thm:covariance}
For every finite $N$ and $\nu\in\cP_2(\R^N)$, the function
$C\mapsto\mathcal F_\nu(C)$ is concave on the positive semidefinite cone.
Equivalently $C\mapsto W_2^2(N(0,C),\nu)$ is convex.
If $\nu$ has $d+1$ affinely independent sites in $\R^d$, each of positive
mass, then $\mathcal F_\nu$ is smooth on the positive definite cone and
its Hessian is negative definite there.
\end{theorem}

Fix $m\geq3$, put $d=m-1$, and identify
$H=\{x\in\R^m:\sum_i x_i=0\}$ isometrically with $\R^d$.
Let
\[
 v_i=e_i-\frac1m\1,\quad
 \mu_m=\E\max_{1\leq i\leq m}Z_i,\quad
 \alpha=\frac{\mu_m}{m-1},\quad
 c_i=\alpha v_i,\quad
 w=\frac{\mu_m^2}{m-1},
\]
where the $Z_i$ are independent standard normals. The regular partition is
$g_i(x)=\1_{\{v_i\cdot x=\max_j v_j\cdot x\}}$, with arbitrary choices on
null tie hyperplanes. Symmetry gives $\E[Xg_i(X)]=c_i$:
the centroids are multiples of $v_i$ and
$\sum_i v_i\cdot\E[Xg_i(X)]=\mu_m$ determines the multiplier.
In larger dimensions the partition ignores the unused coordinates.
With the noise notation below, every benchmark $S_\rho(g)=B_\rho(g,g)$ is
evaluated in dimension $d=m-1$ and is unchanged by adding unused Gaussian
coordinates. When $N<d$, the comparison is with this $d$-dimensional benchmark,
not with an assignment $g$ in $\R^N$.

A \emph{relaxed assignment} is a measurable map into
$\Delta_m=\{t\in[0,1]^m:\sum_i t_i=1\}$. It is balanced if
$\E f_i=1/m$. Write $b_i(f)=\E[Xf_i(X)]$ and
$W(f)=\sum_i\norm{b_i(f)}^2$. Let $\cO_m$ denote the ordered regular
centroid orbit under orthogonal maps. Label permutations are already
implemented by orthogonal symmetries of the simplex.

\begin{theorem}[Balanced endpoint and coercivity]\label{thm:endpoint}
For every balanced relaxed assignment in any finite dimension,
$W(f)\leq w$. If $N\geq d$, equality holds exactly for regular simplex
partitions, up to orthogonal maps, labels, and null changes. There is a finite
$K_m\geq1$, depending only on $m$, such that for all $N\geq d$,
\begin{equation}\label{eq:Hm}
 \dist_{\HS}(b(f),\cO_m)^2\leq K_m\bigl(w-W(f)\bigr).
\end{equation}
For $N<d$ the endpoint benchmark is strict.
\end{theorem}

For $|\rho|<1$ let
$T_\rho u(x)=\E u(\rho x+\sqrt{1-\rho^2}Z)$ and define
\[
 B_\rho(f,h)=\sum_i\ip{f_i}{T_\rho h_i}_{L^2(\gamma_N)},
 \qquad S_\rho(f)=B_\rho(f,f).
\]

\begin{theorem}[Near independence]\label{thm:noise}
For each fixed $m\geq3$ there are numbers $\rho_m^+>0$ and $r_m^->0$,
independent of ambient dimension, with the following properties in every
finite ambient dimension $N$.
\begin{enumerate}[label=(\roman*)]
\item If $0<\rho\leq\rho_m^+$ and $f$ is balanced, then
$S_\rho(f)\leq S_\rho(g)$.
\item If $0<r\leq r_m^-$ and $\E f=\E h=p$ for an arbitrary common
probability vector $p$, then $B_{-r}(f,h)\geq B_{-r}(g,g)$.
\end{enumerate}
For $N\geq d$ equality in (i) has the regular-simplex classification.
Equality in (ii) forces $p=(1/m,\ldots,1/m)$ and $f=h$ equal to one common
regular simplex partition. For $N<d$ both benchmarks are strict.
All statements include relaxed assignments.
\end{theorem}

In particular the negative single-assignment minimum has no imposed mass
constraint. This is not a theorem optimizing a prescribed unequal mass vector.
The two intervals can differ and may shrink with $m$; no numerical bound
uniform in $m$ is claimed. The covariance proof and endpoint argument are
given in the main text. Appendices~\ref{app:common}--\ref{app:negative} give
the complete fixed-dimensional transfers. Section~\ref{sec:dimension} states
the precise external dimension reductions and completes the ambient argument.

\paragraph{Related problems.}
Heilman \cite[Section 8, Conjecture 3]{H2014} asks whether balanced
four-cell first-moment maximizers in $\R^3$ are simplicial conical partitions,
and identifies quantitative stability as a route to larger-cell
near-zero noise results. Theorem~\ref{thm:endpoint} supplies this structure
and stability, and extends them to every $m\geq3$.
The introduction following Theorem 1.3 of Heilman's 23 October 2026
preprint \cite{H2026} still describes the four-cell and higher
Standard Simplex cases as open. That discussion distinguishes the
unrestricted Propeller problem, which does not prescribe cell masses,
from the Standard Simplex problem. It also distinguishes Gaussian
mean-width comparisons with prescribed individual variances.
Our endpoint prescribes cell masses and permits unequal centroid lengths;
it does not assume those variance constraints.
These comparisons locate the related problems and do not constitute an
exhaustive priority assessment. The full-correlation ranges remain open.

